% ---------------------------------------------------------------------------
%  Visual identity for the report. Compile with LuaLaTeX (make, or latexmk).
%  Writers should not need to touch this file.
% ---------------------------------------------------------------------------

% Page ---------------------------------------------------------------------
\usepackage[letterpaper, margin=1in, headheight=14pt, headsep=18pt]{geometry}
\linespread{1.06}
\setlength{\parindent}{1.2em}
\setlength{\parskip}{0pt}
\raggedbottom

% Type: one family across text, headings, figures and math -------------------
\usepackage{fontspec}
\usepackage{unicode-math}
\setmainfont{Libertinus Serif}[Numbers={Proportional,Lining}]
\setsansfont{Libertinus Sans}
\setmonofont{Libertinus Mono}[Scale=MatchLowercase]
\setmathfont{Libertinus Math}
\newfontfamily\displayfont{Libertinus Serif Display}
\usepackage[final, protrusion=true, expansion=true, tracking=true]{microtype}
\usepackage[english]{babel}

% Colour: an accent and its counterpart, as in the figures ------------------
\usepackage{xcolor}
\definecolor{accent}{HTML}{1F6FA8}
\definecolor{brick}{HTML}{B03E26}
\definecolor{ink}{HTML}{15202A}
\definecolor{ink2}{HTML}{4E6273}
\definecolor{rulecolor}{HTML}{C9D4DC}
\definecolor{wash}{HTML}{F3F6F8}
\color{ink}

% Headings -------------------------------------------------------------------
\usepackage[explicit]{titlesec}
\titleformat{\section}
  {\sffamily\fontsize{14}{17}\selectfont\bfseries}
  {\textcolor{accent}{\thesection}}{0.7em}{#1}
\titleformat{\subsection}
  {\sffamily\fontsize{11.5}{14}\selectfont\bfseries}
  {\textcolor{accent}{\thesubsection}}{0.7em}{#1}
\titleformat{\subsubsection}[runin]
  {\sffamily\normalsize\bfseries}{}{0pt}{#1.}
\titlespacing*{\section}{0pt}{1.7em plus 0.4em minus 0.2em}{0.55em}
\titlespacing*{\subsection}{0pt}{1.2em plus 0.3em minus 0.2em}{0.4em}
\titlespacing*{\subsubsection}{0pt}{0.7em plus 0.2em}{0.6em}

% Running header and footer --------------------------------------------------
\usepackage{fancyhdr}
\pagestyle{fancy}
\fancyhf{}
\renewcommand{\headrulewidth}{0.4pt}
\renewcommand{\headrule}{\color{rulecolor}\hrule height\headrulewidth}
\fancyhead[L]{\sffamily\footnotesize\color{ink2}From Listings to Intelligence\enspace·\enspace Midterm report}
\fancyhead[R]{\sffamily\footnotesize\color{ink2}\thepage}
\fancypagestyle{plain}{\fancyhf{}\renewcommand{\headrulewidth}{0pt}%
  \fancyfoot[C]{\sffamily\footnotesize\color{ink2}\thepage}}

% Figures, tables, captions -------------------------------------------------
\usepackage{graphicx}
\usepackage{tikz}
\usetikzlibrary{arrows.meta, calc, positioning}
\graphicspath{{generated/figures/}}
\usepackage{booktabs}
\setlength{\aboverulesep}{0.4ex}
\setlength{\belowrulesep}{0.65ex}
\usepackage{siunitx}
\sisetup{group-separator={,}, group-minimum-digits=4, detect-all}
\usepackage{caption}
\DeclareCaptionLabelFormat{accent}{\textcolor{accent}{#1~#2}}
\captionsetup{font={small,sf}, labelfont=bf, labelformat=accent, labelsep=period,
  justification=justified, singlelinecheck=false, skip=7pt}
\captionsetup[table]{position=top, skip=5pt}
\usepackage{subcaption}
\usepackage{float}
% Floats never appear before the text that introduces them (flafter), and never
% drift past the next section heading (placeins).
\usepackage{flafter}
\usepackage[section]{placeins}
% Pages holding only floats start at the top rather than centring them.
\makeatletter
\setlength{\@fptop}{0pt}
\setlength{\@fpsep}{14pt plus 1fil}
\setlength{\@fpbot}{0pt plus 1fil}
\makeatother
\setlength{\textfloatsep}{14pt plus 3pt minus 3pt}
\setlength{\floatsep}{12pt plus 2pt minus 2pt}
\setlength{\intextsep}{12pt plus 2pt minus 2pt}

% Lists, notes ---------------------------------------------------------------
\usepackage{enumitem}
\setlist{itemsep=0.15em, topsep=0.35em, leftmargin=1.3em}
\usepackage[bottom, hang]{footmisc}
\setlength{\footnotemargin}{0.9em}

% Citations and cross-references --------------------------------------------
\usepackage[style=numeric-comp, sorting=none, maxbibnames=4, giveninits=true,
  url=false, doi=true, backend=biber]{biblatex}
\addbibresource{refs.bib}
\renewcommand*{\bibfont}{\small}
\usepackage[colorlinks, linkcolor=accent, citecolor=accent, urlcolor=accent,
  pdftitle={From Listings to Intelligence: Midterm Report},
  pdfauthor={Columbia Data Science Capstone}]{hyperref}
\usepackage[capitalise, noabbrev]{cleveref}

% Title block: on page 1, no separate title page (it would cost one of ten) --
\newcommand{\titleblock}[5]{%
  \thispagestyle{plain}%
  \begingroup\raggedright
  {\sffamily\footnotesize\bfseries\color{accent}\addfontfeature{LetterSpace=8}%
    \MakeUppercase{#1}\par}
  \vspace{0.7em}
  {\displayfont\fontsize{26}{30}\selectfont #2\par}
  \vspace{0.35em}
  {\fontsize{13.5}{17}\selectfont\itshape\color{ink2} #3\par}
  \vspace{1.0em}
  {\sffamily\small #4\par}
  \vspace{0.25em}
  {\sffamily\footnotesize\color{ink2} #5\par}
  \vspace{0.9em}
  {\color{rulecolor}\rule{\linewidth}{0.5pt}\par}
  \endgroup
  \vspace{0.6em}}

\newcommand{\member}[2]{\mbox{#1\,{\color{ink2}\scriptsize #2}}}

\renewenvironment{abstract}
  {\par\noindent\small\textsf{\textbf{\color{accent}Abstract.}}\enspace\ignorespaces}
  {\par\vspace{0.4em}{\color{rulecolor}\noindent\rule{\linewidth}{0.5pt}}\par\vspace{0.2em}}

% Drafting aids. \draftfalse (or `make final`) removes every brief and marker.
\newif\ifdraft \drafttrue
\usepackage[most]{tcolorbox}
\newtcolorbox{briefbox}[1]{enhanced, breakable, sharp corners, boxrule=0pt,
  colback=wash, left=7pt, right=7pt, top=5pt, bottom=5pt, before skip=6pt, after skip=8pt,
  fontupper=\sffamily\footnotesize\color{ink2}, title={#1}, coltitle=ink2,
  fonttitle=\sffamily\scriptsize\bfseries\addfontfeature{LetterSpace=6}, colbacktitle=wash,
  attach title to upper, after title={\par\smallskip}}
% \brief{owner}{pages}{what to cover}
\newcommand{\brief}[3]{\ifdraft\begin{briefbox}{\MakeUppercase{Writing brief · #1 · #2}}#3\end{briefbox}\fi}
% \figslot{key}{what the figure shows}{height}{note}: holds a figure's place, label fig:key
\newcommand{\figslot}[4]{%
  \begin{figure}[htbp]\centering
  \begin{tcolorbox}[enhanced, sharp corners, colback=white, boxrule=0pt, height=#3, valign=center,
    halign=center, borderline={0.6pt}{0pt}{rulecolor, dashed}]
  {\sffamily\footnotesize\color{ink2}\textbf{Figure to come.}\enspace #4}
  \end{tcolorbox}
  \caption{\textcolor{ink2}{Placeholder: #2.}}\label{fig:#1}\end{figure}}
\newcommand{\todo}[1]{\ifdraft\textcolor{brick}{\sffamily\footnotesize[#1]}\fi}
